% ============================================================================
%  金港城物业费小程序 · 使用手册（业主篇 / 员工篇）
%
%  编译： xelatex manual.tex        （跑两遍，第二遍生成目录页码）
%  依赖： TeXLive 2023+ 的 ctex / tcolorbox / titlesec / fontspec
%  插图： img/ 目录；换图只需同名覆盖，无需改本文件
% ============================================================================
\documentclass[11pt,a4paper]{article}

\usepackage[UTF8,fontset=macnew]{ctex}
\usepackage[margin=2.1cm,top=2.3cm,bottom=2.2cm,headsep=0.7cm]{geometry}
\usepackage{graphicx}
\usepackage{xcolor}
\usepackage{tcolorbox}
\usepackage{titlesec}
\usepackage{fancyhdr}
\usepackage{array}
\usepackage{enumitem}
\usepackage{hyperref}
\usepackage{tikz}
\tcbuselibrary{skins}

% ── 取小程序自己的配色，手册和界面是同一家人 ───────────────────────────────
\definecolor{plum}{HTML}{3B2B57}      % 主色：深紫
\definecolor{deepink}{HTML}{251C38}
\definecolor{gold}{HTML}{C9A66B}      % 强调：金
\definecolor{golddeep}{HTML}{B8862F}
\definecolor{cream}{HTML}{FBF7F0}     % 底：暖白
\definecolor{cream2}{HTML}{F3EDE3}
\definecolor{inkgray}{HTML}{55496B}
\definecolor{subgray}{HTML}{7C7288}
\definecolor{danger}{HTML}{C45656}
\definecolor{okgreen}{HTML}{4F7D44}
\definecolor{rule}{HTML}{E2DDD4}

\hypersetup{colorlinks=true,linkcolor=plum,urlcolor=golddeep,pdftitle={金港城物业费小程序使用手册},pdfauthor={港城物业}}

% ── 版式 ─────────────────────────────────────────────────────────────────
\setlength{\parindent}{0pt}
\setlength{\parskip}{0.35em}
\linespread{1.28}

\titleformat{\section}[block]
  {\normalfont\LARGE\bfseries\color{plum}}{\thesection}{0.6em}{}
\titleformat{\subsection}[block]
  {\normalfont\large\bfseries\color{deepink}}{}{0em}{}
\titlespacing{\section}{0pt}{1.4em}{0.8em}
\titlespacing{\subsection}{0pt}{1.1em}{0.4em}

\setlength{\headheight}{22pt}
\pagestyle{fancy}
\fancyhf{}
\renewcommand{\headrulewidth}{0.4pt}
\renewcommand{\headrule}{\hbox to\headwidth{\color{rule}\leaders\hrule height \headrulewidth\hfill}}
\fancyhead[L]{\small\color{subgray}金港城物业费小程序 · 使用手册}
\fancyhead[R]{\small\color{subgray}\leftmark}
\fancyfoot[C]{\small\color{subgray}\thepage}

% ── 步骤块：左图右文，比例固定，全篇统一 ──────────────────────────────────
%   \step{图片文件名}{步骤标题}{正文}
%   图片按高度对齐（手机截图都是竖的），所以每一步的视觉重量一致。
% 插图全部预处理成 375×760(见 scratchpad/manual-shots.mjs 与裁切脚本),
% 所以这里按**宽度**定尺寸即可 —— 每一步的图大小完全一致，版面才真的统一。
% 早先按高度对齐:手机截图太瘦长，被压成一条细缝，字根本看不清。
\newlength{\shotW}\setlength{\shotW}{5.15cm}  % 左栏宽度（含边框留白）
\newlength{\shotIW}\setlength{\shotIW}{4.85cm} % 图片净宽
\newcounter{stepno}[section]

\newtcolorbox{shotframe}{
  enhanced, sharp corners=downhill, arc=3pt, boxrule=0.5pt,
  colback=cream, colframe=rule, boxsep=0pt, left=2pt, right=2pt, top=2pt, bottom=2pt,
  drop shadow={black!12}
}

\newcommand{\stepbadge}[1]{%
  \begin{tikzpicture}[baseline=(n.base)]
    \node[circle, fill=plum, text=cream, font=\bfseries\small, inner sep=2.4pt, minimum size=6.4mm] (n) {#1};
  \end{tikzpicture}%
}

\newcommand{\step}[3]{%
  \par\vspace{0.55em}\noindent
  \refstepcounter{stepno}%
  \begin{minipage}[t]{\shotW}
    \vspace{0pt}
    \begin{shotframe}
      \includegraphics[width=\shotIW]{img/#1}
    \end{shotframe}
  \end{minipage}%
  \hfill
  \begin{minipage}[t]{\dimexpr\textwidth-\shotW-0.7cm\relax}
    \vspace{0pt}
    {\stepbadge{\thestepno}\hspace{0.35em}\textbf{\large\color{deepink}#2}}\par
    \vspace{0.45em}
    #3
  \end{minipage}%
  \par\vspace{1.0em}
}

% ── 三种提示框：为什么这么设计 / 注意 / 这是钱 ─────────────────────────────
% 注意:这两个框会用在 \step 的 minipage 里,那里**不能** breakable ——
% tcolorbox 会把框拆成多段而正文一段都放不下,页面上只剩几条空色条(2026-08-05 修)。
\newtcolorbox{why}{
  enhanced, colback=cream2, colframe=gold, boxrule=0pt,
  leftrule=2.4pt, arc=2pt, left=8pt, right=8pt, top=5pt, bottom=5pt,
  fontupper=\small\color{inkgray}
}
\newtcolorbox{warn}{
  enhanced, colback=danger!5, colframe=danger, boxrule=0pt,
  leftrule=2.4pt, arc=2pt, left=8pt, right=8pt, top=5pt, bottom=5pt,
  fontupper=\small\color{inkgray}
}
\newcommand{\kw}[1]{\textbf{\color{plum}#1}}
% 行内的「按钮」标签:内边距压到最小 —— 否则一行里出现一个按钮就把整段行距顶开
\newcommand{\btn}[1]{\tcbox[on line, colback=plum, colframe=plum, boxsep=0pt,
  left=2.5pt, right=2.5pt, top=0.6pt, bottom=0.6pt, arc=1.6pt]{\color{cream}\footnotesize\textbf{#1}}}
\newcommand{\gbtn}[1]{\tcbox[on line, colback=gold!30, colframe=gold, boxsep=0pt,
  left=2.5pt, right=2.5pt, top=0.6pt, bottom=0.6pt, arc=1.6pt]{\color{deepink}\footnotesize\textbf{#1}}}

\begin{document}

% ══════════════════════════════ 封面 ══════════════════════════════
\begin{titlepage}
\thispagestyle{empty}
\begin{tikzpicture}[remember picture, overlay]
  \fill[plum] (current page.north west) rectangle ([yshift=-11.5cm]current page.north east);
  \fill[gold] ([yshift=-11.5cm]current page.north west) rectangle ([yshift=-11.72cm]current page.north east);
\end{tikzpicture}

\vspace*{2.6cm}
{\color{cream}\fontsize{31}{38}\selectfont\bfseries 物业费小程序}\par
\vspace{0.5em}
{\color{gold}\fontsize{31}{38}\selectfont\bfseries 使用手册}\par
\vspace{1.6em}
{\color{cream!85}\large 业主篇 · 物业员工篇}\par
\vspace{0.7em}
{\color{cream!60}\normalsize 交物业费、看账单、报维修；发账单、收现金、催缴、管员工}\par

% 版权行压到页脚 —— 原来靠 \vspace{6.6cm} 顶下去,是因为下面还有一段说明文字;
% 那段删了之后固定值就把它吊在半空,改成 \vfill 让它自己贴底。
\vfill
{\color{inkgray}\large 金港城 · 港城物业}\par
\vspace{0.3em}
{\color{subgray}\small 适用版本 1.4.0 \quad|\quad 2026 年 8 月}\par
\vspace{1.4cm}
\end{titlepage}

% ══════════════════════════════ 目录 ══════════════════════════════
\thispagestyle{empty}
{\color{plum}\Large\bfseries 目录}\par\vspace{0.6em}
{\color{rule}\hrule height 0.6pt}\vspace{0.9em}
\begingroup
\hypersetup{linkcolor=deepink}
\makeatletter
\@starttoc{toc}
\makeatother
\endgroup
\newpage

% ══════════════════════════════ 业主篇 ══════════════════════════════
\section{业主篇：交费只有一件事}
\markboth{业主篇}{业主篇}

\begin{why}
业主要做的事只有一件：\kw{看到自己家的账单，把它交掉}。所以整个流程被压到
「授权手机号 → 看账单 → 缴费」三步，中间不需要注册、不需要记密码、不需要填房号 ——
房号是物业登记好的，系统按手机号自动对上。
\end{why}

\step{o-mine-admin}{第一次进来：授权手机号}{%
打开小程序，底部点\kw{「我的」}，按提示\kw{授权手机号}。
这是唯一一次需要你点确认的地方 —— 微信会问「是否允许获取你的手机号」，点允许。

授权之后，页面上会出现\kw{「我的房屋」}，显示你家的小区与房号。
这说明物业已经把这个号码登记到你家名下了，账单会自动出现，不用再做任何事。
}

\step{o-bind}{如果提示「物业未登记此号码」}{%
说明物业那边登记的是另一个号（比如买房时留的号）。两种办法：

\begin{enumerate}[leftmargin=1.4em, itemsep=0.2em]
  \item \kw{告诉物业你现在的号码}，他们在后台加一下，你重新进小程序就能看到账单 —— 最快；
  \item 自己\kw{申请绑定}：选小区 → 输房号 → 填姓名 → 提交，等物业审核通过。
\end{enumerate}

房号支持多种写法，\mbox{「1栋101」}\mbox{「1-101」}\mbox{「101」}都能搜到，不用纠结格式。

\begin{why}
为什么手机号能直接对上房屋：物业在系统里给每套房登记了\kw{授权手机号}。
凡是在名单里的号码，微信授权后立刻看到那套房的账单；一旦物业把号码移除，
你也立刻看不到 —— 这样换租、换房主都不用等谁去解绑。
\end{why}
}

\step{o-bill}{看账单：按年度分组，一眼看出还欠多少}{%
底部\kw{「账单」}进来。最上面深色卡片里的数字是\kw{待缴合计}，
它取自系统的权威汇总，不受你当前翻到第几页影响。

下面按\kw{年度}分组。每一笔写清费用名称、到期日、金额；
待缴的那笔右边有金色的 \gbtn{缴费} 按钮，已缴的显示缴费日期。

上方三个标签\kw{全部 / 待缴 / 已缴}用来筛选；如果物业收多种费用（物业费、车位费…），
还会多出一行按费用种类的筛选。
}

\step{o-bill-detail}{点开一笔，能看到钱是怎么算出来的}{%
点账单任意一行进入详情。这里给出\kw{完整的计算依据}：

\begin{itemize}[leftmargin=1.4em, itemsep=0.15em]
  \item 计费方式（按建筑面积 / 固定金额）
  \item 建筑面积、单价、月数
  \item 账期起止（\texttt{2026-08-05 \textasciitilde{} 2027-08-04}）
\end{itemize}

\begin{why}
为什么把算式摊开：物业费最常见的纠纷是「这个数怎么来的」。
面积 × 单价 × 月数三个数都摆在这里，你可以拿着房本自己核一遍。
账期按\kw{你家的放户日期}起算，所以邻居的账期和你不一定一样 —— 这是正常的。
\end{why}
}

\step{o-ticket}{报事报修：写清楚哪里坏了}{%
「我的 → 我的报修」或首页的报修入口进来。选\kw{类型}（报修 / 投诉 / 建议），
写清楚\kw{哪里、什么问题}，可以拍照上传。

提交后能看到处理进度：\kw{待受理 → 处理中 → 已办结}。
物业受理时会填处理人姓名，办结时会写回复 —— 这两样你都看得到，还能给这次服务评分。
}

\begin{warn}
\textbf{缴费提醒收不到？}\quad 微信的订阅消息是\kw{一次性}的：每授权一次，才能收到一条。
在「我的 → 缴费提醒」里点开启，可以多授权几次。没授权不影响你看账单和缴费，
只是账单生成、到期前的提醒推不到你微信里。
\end{warn}

\newpage
% ══════════════════════════════ 员工篇 ══════════════════════════════
\section{员工篇：手机就是收费台}
\markboth{员工篇}{员工篇}

\begin{why}
物业这一端全部做在\kw{手机上}，不需要电脑。判断你是不是员工靠\kw{手机号}：
号码在管理员名单里，「我的」页底部就会多出\kw{「物业工作」}这一段；
普通业主看不到它的任何痕迹。所以员工和业主用的是\kw{同一个小程序}，不用装两个。
\end{why}

\subsection{一、进入与总览}

\step{s-entry}{进入管理端}{%
用\kw{已登记为员工的手机号}的微信打开小程序 → 「我的」→ 最下面
\kw{「物业工作 · 物业管理」}。

这一段下面写着你的\kw{角色}：\kw{物业管理员}或\kw{收费员}。
角色决定你能做什么，见本篇最后一节。

\begin{warn}
进不去、显示「没有管理权限」？先确认两件事：这个微信授权的是不是登记过的号码；
这个员工账号是不是\kw{在职}状态。确认后重新打开小程序、重新授权手机号。
\end{warn}
}

\step{s-grid}{楼盘图：颜色就是状态}{%
进来第一屏是\kw{楼盘图}。上面一排是楼栋，点一下切换；下面按\kw{层}铺开房间，
\kw{顶层在上}，和你站在楼前看到的一样。

\begin{itemize}[leftmargin=1.4em, itemsep=0.15em]
  \item 灰白格 = 正常；\textcolor{danger}{\textbf{红格}} = 有待缴，格子上直接写欠多少
  \item 虚线空格 = 已停用的房屋（不出账、不进欠费）
\end{itemize}

顶上五个标签：\kw{楼盘图 / 欠费 / 报修 / 待发布 / 功能}，
数字角标告诉你哪一块有事。\kw{「功能」}是工具箱 —— 发账单、新增房屋、绑定审批、
核销卡券、发公告、员工与权限都在里面。最上面的搜索框用来\kw{接电话查户}：
房号、业主姓名、手机号都能搜。
}

\step{s-tools}{「功能」页签：所有工具都在这里}{%
第五个标签\kw{「功能」}是工具箱，不用再到处找按钮：

\vspace{0.3em}
\begin{tabular}{@{}>{\bfseries\color{plum}}l@{\hspace{0.7em}}p{\dimexpr\linewidth-7.2em\relax}@{}}
发账单 & 批量出账，先草稿后发布\\[0.2em]
新增房屋 & 建房、登记住户手机、挂收费标准\\[0.2em]
绑定审批 & 业主自助申请绑房的审核（有待审时带数字角标）\\[0.2em]
核销卡券 & 扫业主的券码，当面兑奖品\\[0.2em]
发公告 & 停水停电、缴费通知，可置顶\\[0.2em]
员工与权限 & 加人、改角色、停用（仅物业管理员可见）\\
\end{tabular}
}

\step{s-arrears}{欠费与催缴：先打电话}{%
切到\kw{「欠费」}标签。最上面是全小区欠费合计与户数，下面逐户列出欠几笔、逾期几天。

\kw{每一行都有「打电话」} —— 催缴最有效的手段是打电话，所以它排在最前面。
需要群发时，勾选若干户，底部\btn{给选中的 N 户发提醒}。

\begin{warn}
群发的是\kw{微信订阅提醒}，业主没授权过就收不到；
而且同一张账单的同一类提醒\kw{系统只发一次}（微信本身也不允许重复推送）。
所以发完如果提示「这次没有新提醒发出」，不是坏了 —— 是这几户已经提醒过了，
接下来只能打电话。
\end{warn}
}

\subsection{二、一户人家的全部事情}

\step{s-house}{点房间 = 进这一户}{%
在楼盘图上点任意一个房间，进入\kw{这一户的作业台}。头卡是这户的身份与欠缴，
下面四个标签把事情分开，\kw{一屏一件事}：

\vspace{0.3em}
\begin{tabular}{@{}>{\bfseries\color{plum}}l@{\hspace{0.8em}}l@{}}
账单 & 每一笔的状态，以及这笔钱该做的动作\\
住户 & 谁能看到这套房的账单（换租就改这里）\\
标准 & 这户按哪条收费标准出账\\
信息 & 面积、放户日期、启用/停用\\
\end{tabular}
\vspace{0.3em}

账单那一栏里，动作跟着状态给：待缴 → \gbtn{收到现金·登记}；
微信已缴 → \btn{退款}；线下已缴 → \btn{冲正线下收款}；
已退款/已作废 → \btn{重开账单}。底部金色按钮是\kw{只给这一户发账单}。
}

\step{s-house-contacts}{换租、换房主：删旧号 + 加新号}{%
切到\kw{「住户」}。这里列的是\kw{能看到这套房账单的手机号}。

换住户就两步，十秒钟：
\begin{enumerate}[leftmargin=1.4em, itemsep=0.2em]
  \item \btn{移除}旧住户的号 —— 他\kw{立刻}看不到这套房的账单；
  \item 填新住户的手机号，\kw{＋ 添加授权手机号} —— 他在小程序授权手机号后立刻能看到。
\end{enumerate}

移除时系统会告诉你「已同时解除 N 人的绑定」，做了什么一目了然。
姓名只是备注，写不写都行。

\begin{why}
为什么不做「过户流程」：现实里房子频繁出租出售，等前一位住户自己去解绑是等不到的。
所以权限完全由\kw{物业手里的名单}决定 —— 名单里有就看得到，删掉就看不到，不需要谁配合。
\end{why}
}

\step{s-house-info}{面积和放户日期：它们直接决定收多少}{%
切到\kw{「信息」}→\kw{编辑这些信息}。四项可改：房号、建筑面积、放户日期、启用/停用。

\begin{itemize}[leftmargin=1.4em, itemsep=0.15em]
  \item \kw{建筑面积} —— 按面积收费时，它乘进金额里
  \item \kw{放户日期} —— 决定这户\kw{每年几月出账}，页面上直接写着「每年 8 月出账」
  \item \kw{停用} —— 停用的房屋不出账、不进欠费、楼盘图里变灰（空置就该这样）
\end{itemize}

改面积或放户日期时会先弹一句话说清后果，保存后留审计痕迹。

\begin{warn}
\textbf{删房屋}在这一屏最下面的红色区域。出过账单的房屋\kw{删不掉} ——
账单是审计凭据，得能回答「这房去年收过多少」。这种情况请改成\kw{停用}。
确实是测试房要清干净，可以升级到「彻底删除」，那一步要求\kw{原样打出房号}才执行。
\end{warn}
}

\subsection{三、出账与收钱}

\step{s-billone}{给一户补账单：零选择}{%
房屋详情底部\gbtn{给这户发账单}进来。这一页\kw{不需要你做任何选择}：

\begin{itemize}[leftmargin=1.4em, itemsep=0.15em]
  \item 收费标准 = 这户挂着的那条
  \item 月份 = 按这户的\kw{放户日}自动推出来（按年收费只在放户那个月出账）
  \item 金额与算式直接摆出来：\texttt{92.08㎡ × 1.4 元 × 12 个月}
\end{itemize}

确认无误 → \btn{生成这户的账单}（这是\kw{草稿}，业主还看不到）→ 再确认一次才\kw{发布}。

\begin{warn}
如果这张账单被放进了一个\kw{含多户的草稿批次}，页面会告诉你「这一批共 N 户」，
并让你去整批核对后一起发 —— 发布是\kw{整批}发布，不会在单户页面上悄悄把几十户一起发出去。
\end{warn}
}

\step{s-billing}{批量出账：默认全出，点掉不该出的}{%
\kw{「功能」}页签 → \btn{发账单}。三步：选\kw{标准 + 月份 + 范围（全部 / 按楼栋）}
→ \kw{预览}→ 生成草稿 → 发布。

预览页逐户列出金额和依据。\kw{默认全部都出}，点一行就是\kw{本次不给这户出}
（那一行变灰并标「本次不出」）。顶上的合计\kw{跟着实时变} ——
剔掉几户，数字立刻减下去。

\begin{why}
为什么要能点掉：每月总有几户线下已经交过了。
剔掉的户会记进「跳过」并写明原因，下个月有人问「这户为什么没账单」，那里就是答案。
\end{why}
}

\step{s-batches}{待发布：发布前最后一次核对}{%
首页\kw{「待发布」}标签进来（有草稿时标签上有数字角标）。这里是所有\kw{草稿批次}
（业主还看不到的账单）。

\begin{itemize}[leftmargin=1.4em, itemsep=0.15em]
  \item \btn{逐户核对} —— 展开这一批的每一户
  \item \kw{剔除} —— 这户本期不发（要填原因：线下已收 / 不该出）
  \item \btn{发布(N 户)} —— 按钮上的户数取自库里的真实条数，不是生成时的数字
  \item \kw{这批整个不发} —— 算错了就整批作废，重新出
\end{itemize}

\begin{warn}
发布是\kw{不可逆}的一步：业主立刻能看到并缴费。所以按钮上永远写着真实户数，
点下去还要再确认一次。
\end{warn}
}

\step{s-collect}{业主拿现金来交：登记收款}{%
房屋详情的待缴账单行 → \gbtn{收到现金·登记}。

金额\kw{不给改} —— 收多少是账单说的。你要填的是：
\kw{收款方式}（现金 / 微信收款码 / 银行转账 / 其他）、\kw{单据号}（收据本编号，必填）、
收款日期、付款人（选填）。

\begin{why}
为什么单据号必填：这是这笔现金\kw{唯一的线下痕迹}。以后对账对不上，
只能靠这个编号回到纸质收据本。系统给了个默认值（如 \texttt{CASH-0805-A-3-702}），
改成你收据本上的真实编号即可。
\end{why}
}

\step{s-collect-done}{收完给凭证}{%
确认后立刻拿到\kw{收据号}和订单号。业主那边这笔账单当场变「已缴」，
并能看到电子收据。

\begin{itemize}[leftmargin=1.4em, itemsep=0.15em]
  \item 手抖点了两下\kw{不会重复收款}（同一张账单 + 同一个单据号算同一次）
  \item 收错了户 → 回那笔账单点\btn{冲正线下收款}（限管理员），账单回到待缴
\end{itemize}
}

\subsection{四、日常事务}

\step{s-tickets}{报事报修：受理和办结都要填字}{%
首页\kw{「报修」}标签。\kw{待受理}排在最上面 —— 业主在等。

\begin{itemize}[leftmargin=1.4em, itemsep=0.15em]
  \item \btn{受理(填处理人)} —— 填谁去处理，\kw{业主会看到这个名字}
  \item \gbtn{办结(填回复)} —— 写清怎么解决的，\kw{业主会看到这段话，还能评分}
\end{itemize}

\begin{why}
为什么没有「一键办结」：一键办结留下的是一条没人看得懂的记录，
业主只会再报一次。填一句话的成本，换的是不用再来第二遍。
\end{why}
}

\step{s-announce}{发公告}{%
\kw{「功能」}页签 → \btn{发公告}。填标题和正文，两个开关：
\kw{发给全部小区}（多小区时用）、\kw{置顶}。

发布前的确认框会写清\kw{受众范围}。下面是已发布列表，每条可\kw{撤回} ——
但撤回只是让它不再显示，\kw{已经看过的人收不回来}。
}

\step{s-coupon}{卡券核销:业主拿券来兑奖品}{%
\kw{「功能」}页签 → \kw{核销卡券}(收费员就能操作)。业主打开「我的 → 我的卡券 →
\kw{亮码核销}」出示二维码,你点 \btn{扫业主的券码} 扫一下 → 看清\kw{是什么券、
能不能用} → \gbtn{确认核销}。业主手机不方便时,手动输券码也一样。

\begin{itemize}[leftmargin=1.4em, itemsep=0.15em]
  \item 核销\kw{不可逆}:确认前把东西(电影票等)准备好当面交付
  \item 同一张券\kw{核不了第二次} —— 两个前台同时核,只有一个成功
  \item 已过期、已用过的券会直接告诉你,不用自己看日期
\end{itemize}
}

\step{s-housenew}{新增房屋}{%
\kw{「功能」}页签 → \kw{新增房屋}。一屏填完：房号、显示名、类型（住宅/门市/车位）、
建筑面积、放户日期、住户手机、收费标准。
房号下面会\kw{实时写出}这套房归入楼盘图的哪栋哪层；认不出的写法会给出参考格式。

\begin{itemize}[leftmargin=1.4em, itemsep=0.15em]
  \item 房号在本小区内\kw{唯一}。填了已有的房号是\kw{更新}那套房，不会出现两套同号
  \item 面积和放户日期可以以后补，但\kw{补齐之前这套房出不了账}
  \item 住宅\kw{必须}填面积（按面积收费算不出金额）
  \item 填了住户手机，对方在小程序授权手机号后立刻能看到这套房的账单
\end{itemize}
}

\subsection{五、员工与权限（仅物业管理员）}

\begin{why}
系统\kw{不允许}你把自己停用或降级，也\kw{不允许}拿掉公司里最后一个在职管理员 ——
否则一次手滑就没人能管员工、退款、发布账单了，恢复得找开发改数据库。
想交接管理员：先把接手的人\kw{设成管理员}，再把自己降成收费员。
\end{why}

\step{s-staff}{加人、改角色、停用}{%
\kw{「功能」}页签 → \kw{员工与权限}（只有\kw{物业管理员}看得到这个入口）。

\vspace{0.3em}
\begin{tabular}{@{}>{\bfseries\color{plum}}l@{\hspace{0.7em}}p{\dimexpr\linewidth-6.2em\relax}@{}}
收费员 & 收钱、催缴、出账、报修；\textbf{退款、整批作废做不了}\\[0.3em]
物业管理员 & 什么都能做：退款、整批作废、彻底删房、管员工\\
\end{tabular}
\vspace{0.7em}

每个人一行，可以改\kw{手机号}、\kw{切换角色}、\kw{重置密码}、\kw{停用/恢复}。

\begin{warn}
\textbf{三件必须知道的事}
\begin{enumerate}[leftmargin=1.4em, itemsep=0.2em, topsep=0.25em]
  \item \kw{初始密码只显示一次}，忘了只能重置。
  \item 填了手机号的员工\kw{授权手机号即可进入}；初始密码只在用电脑后台时需要。
  \item \kw{停用和改角色即时生效}，他手里的登录状态当场失效。
\end{enumerate}
\end{warn}
}

\end{document}
